%! Author = sriadityavedantam
%! Date = 3/30/26

\section{Homomorphisms and Isomorphisms}

\begin{definition}[Homomorphism]
    Let $\ri$ and $\si$ be rings.
    Suppose a function $f:\ri\to\si$ is given.
    If for all $a,b\in\ri$,
    \[
        f(a+b)=f(a)+f(b)
    \]
    and
    \[
        f(ab)=f(a)f(b),
    \]
    then $f$ is called a homomorphism.
\end{definition}

\begin{definition}[Isomorphism]
    Suppose $f:\ri\to\si$ is a bijective homomorphism.
    Then $f$ is called an isomorphism.
\end{definition}

If $f$ is an isomorphism, then $\ri$ and $\si$ are isomorphic to each other.
Suppose $\ri=\mathbb{Z}$ and $\si=2\mathbb{Z}$, and let $f:\ri\to\si$ be defined by $f(m)=2m$.
Is this an isomorphism?

\begin{example}
    We compute
    \[
        f(m+n)=2(m+n)=2m+2n=f(m)+f(n).
    \]
    So $f$ preserves addition.
    However,
    \[
        f(mn)=2mn,
    \]
    while
    \[
        f(m)f(n)=(2m)(2n)=4mn.
    \]
    Thus
    \[
        f(mn)\neq f(m)f(n)
    \]
    in general.
    So $f$ is not a ring homomorphism.
    Hence $\mathbb{Z}$ and $2\mathbb{Z}$ are not isomorphic by this map.
\end{example}

\begin{definition}[Endomorphism]
    A homomorphism that maps $\ri$ to itself is called an endomorphism.
\end{definition}

\begin{definition}[Automorphism]
    An isomorphic endomorphism is called an automorphism.
\end{definition}

\begin{proposition}{
    A function is bijective if and only if it has an inverse.
}
    First suppose $f:\ri\to\si$ has an inverse function $g:\si\to\ri$.
    Then
    \[
        g\circ f=\id_{\ri}
    \]
    and
    \[
        f\circ g=\id_{\si}.
    \]
    Hence $f$ is injective and surjective, so $f$ is bijective.
    Conversely, suppose $f$ is bijective.
    Since $f$ is surjective, for every $y\in\si$ there exists $x\in\ri$ such that $f(x)=y$.
    Since $f$ is injective, this $x$ is unique.
    So we may define a function $g:\si\to\ri$ by letting $g(y)=x$, where $f(x)=y$.
    Then
    \[
        g\circ f=\id_{\ri}
    \]
    and
    \[
        f\circ g=\id_{\si}.
    \]
    Thus $g$ is an inverse of $f$.
\end{proposition}

\begin{axiom}[Isomorphism Properties]
    We can check some properties of a homomorphism to rule out whether two rings are isomorphic.
    In particular, isomorphic rings must have the same:
    \begin{enumerate}
        \item number of elements, when finite;
        \item number of units;
        \item number of zero-divisors.
    \end{enumerate}
    These are useful necessary checks, though by themselves they do not always prove two rings are isomorphic.
\end{axiom}

For example, we can state that $\z\not\cong\q$.
The reason is that every nonzero element in $\q$ is a unit, while the only units in $\z$ are $\pm 1$.
Similarly, $\z_4\not\cong\z_6$, due to the number of elements.

Perhaps in previous courses, such as Calculus III, you have looked at $\r^3$, which means a $3$-tuple ordered triple that represents $(x,y,z)$ in space.
However, this is a generalized fact.
What if I wanted to have two points from different sets, but still create an ordered pair or tuple?

\begin{axiom}[Cartesian Product]
    If $\ri$ and $\si$ are rings, then
    \[
        \ri\times\si\coloneqq\{(r,s):r\in\ri,\ s\in\si\}
    \]
    is also a ring under addition and multiplication defined by
    \begin{align*}
        (r_{1},s_{1})+(r_2,s_2)&=(r_1+r_2,s_1+s_2),\\
    (r_1,s_1)(r_2,s_2)&=(r_{1}r_2,s_{1}s_2).
    \end{align*}
\end{axiom}

It will be a fun exercise to prove the following lemma, or at least work through a couple of examples.

\begin{lemma}{
    If $\gcd(m,n)=1$, then
    \[
        \z_m\times\z_n\cong\z_{mn}.
    \]
}
\end{lemma}

Note that in $\zxz$, the zero-divisors include $(0,1)$ and $(1,0)$.
Let $\ri=\si=\z$ in $\zxz$.
Let $\pi:\zxz\to\z$ be the projection map onto the first coordinate.
Then
\[
    \pi((1,0))=1,
\]
which is a unit.
So a homomorphism need not preserve zero-divisors.